% =============================================================================
% Anhang A: Provenienzregister
% =============================================================================

\chapter{Provenienzregister}
\label{app:provenienz}

Dieses Register setzt das in \secref{sec:provenienz} beschriebene Schema um. Es wird fortlaufend und
rückwirkend gefüllt; Lücken tragen die Werte \zitat{unbekannt} oder \zitat{nicht protokolliert}. Interne
Belege (Freeze-Dossiers, Sitzungsprotokolle, Kampagnenverzeichnisse) liegen aus Datenschutzgründen
unversioniert beim Autor; Archive in \file{thesis/evidence/} sind über ihre SHA-256-Prüfsummen identifizierbar.

\begin{footnotesize}
\begin{longtable}{@{}L{3.7cm}L{1.7cm}L{2.3cm}L{1.5cm}L{2.0cm}L{1.3cm}@{}}
  \caption{Provenienzregister (Arbeitsstand 2026-08-13)}
  \label{tab:provenienzregister} \\
  \toprule
  \textbf{Aussage / Artefakt} & \textbf{Erzeugung} & \textbf{unabh. Prüfung} & \textbf{Datum} &
    \textbf{Beleg} & \textbf{Status} \\
  \midrule
  \endfirsthead
  \toprule
  \textbf{Aussage / Artefakt} & \textbf{Erzeugung} & \textbf{unabh. Prüfung} & \textbf{Datum} &
    \textbf{Beleg} & \textbf{Status} \\
  \midrule
  \endhead
  \bottomrule
  \endfoot
  Zählung: 103 normative \texttt{>{}>}-Absätze in RFC~9868
    & Werkzeug und Modell
    & unabhängige Nachzählung mit zweitem Verfahren; früherer Wert 96 verworfen
    & 2026-08-13 (letzte Prüfung)
    & \path{literature/rfc9868.txt}
    & bestätigt \\
  drei handgerechnete OCS-Referenzvektoren, nicht palindromisch
    & Autor und Modell
    & unabhängige Nachrechnung; externe Bestätigung offen
    & 2026-08-13
    & Impl.-Repo, Commit \texttt{6fa4ca1}
    & offen \\
  Kampagnenarchiv\newline\path{external-campaign-}\newline\path{20260810T200118Z}\newline\path{.tar.zst}
    & Messwerkzeuge
    & SHA-256 gegen \file{.sha256}-Datei
    & 2026-08-13 (Prüfung)
    & \path{thesis/evidence/}
    & bestätigt \\
  Kampagnenarchiv\newline\path{bidir-campaign-}\newline\path{20260811.tar.zst}
    & Messwerkzeuge
    & SHA-256 gegen \file{.sha256}-Datei; Prüfsummen dreifach nachgerechnet
    & 2026-08-13 (Prüfung)
    & \path{thesis/evidence/}
    & bestätigt \\
  GitHub-Zustand beider Repositories (Rulesets, Workflows, kein verpflichtender Check)
    & Werkzeug (API-Abfrage)
    & zweite Abfrage an zwei Tagen
    & 2026-08-12 und 2026-08-13
    & Freeze-Dossier (lokal, unversioniert)
    & bestätigt \\
  Audit-Bericht des Zweitmodells zum Implementierungsstand
    & Zweitmodell
    & Wiederherstellung aus Sitzungsprotokoll nach Bedienfehler
    & 2026-07-10 bis 2026-08-13
    & Impl.-Repo, \file{review.md} (unversioniert)
    & vorhanden \\
  Kapitelstart-Entscheidung (Methodik zuerst)
    & Autor und führendes Modell
    & ein unabhängiges Claude- und zwei Codex-Gutachten, nur lesend
    & 2026-08-13
    & Sitzungsprotokoll (lokal)
    & bestätigt \\
  Kapitel-3-Entwurf (Methodik)
    & führendes Modell
    & adversariale Prüfung durch das Zweitmodell; Befunde eingearbeitet
    & 2026-08-13
    & Sitzungsprotokoll (lokal)
    & bestätigt \\
  Quellenprüfung der externen Referenzfälle (\secref{sec:externe-faelle})
    & Werkzeug und Modell
    & Autor-Hinweis deckte einen vom Abrufwerkzeug ausgelassenen Abstract-Satz auf; Erstbewertung korrigiert
    & 2026-08-13
    & Sitzungsprotokoll (lokal)
    & bestätigt \\
  Kampagnen-Kontrollzahl in Abschnitt~3.4 (Notiz 15, belegt 18)
    & führendes Modell
    & Zweitmodell-Befund; eigene Nachprüfung am Kampagnenbericht (Erratum E1)
    & 2026-08-13
    & Impl.-Repo, Kampagnenbericht
    & bestätigt \\
  Modell- und Konfigurationsprovenienz früher Projektphasen
    & unbekannt
    & entfällt
    & nicht protokolliert
    & entfällt
    & offen \\
\end{longtable}
\end{footnotesize}
